% ===========================================================================
% Compact 2D Gaussian Captures for Memory-Efficient 3D Gaussian Splatting
% Working draft — two claim-separated parts:
%   Part I  : the systems / VRAM claim (image-free compact-field supervision)
%   Part II : Beam Fusion (tomographic initialization) + carrier refinement
%
% Draft conventions: BLACK text is explained and evidence-bound; RED text
% (and red framed figure slots) still has to be shown/measured/frozen.
% See preamble.tex for the macros and paper/README.md for the build.
%
% Source of truth for the claim discipline:
%   docs/PAPER_PLAN_beam_fusion.md, docs/20260725_claims_and_questions.md,
%   docs/THREE_ARM_EXPERIMENT_PROGRAM.md, ara/logic/claims.md
% ===========================================================================
\documentclass[11pt,a4paper]{article}

% ---------------------------------------------------------------------------
% Preamble and draft macros.
%
% Draft conventions (owner decision, 2026-07-30):
%   * Black text is explained and evidence-bound. Statements backed by
%     repository artifacts carry gray [evidence ...] notes.
%   * Red text is a TODO. Every red item starts with "TODO:". Red framed
%     boxes inside figure floats are figure slots that specify which figure
%     has to be produced and from which data.
%   * Style rules for the running text: no em dashes, no semicolons, no
%     colons (the "TODO:" prefix is the one exception), no marketing
%     vocabulary. Negative results appear only where an obvious approach
%     does not apply and an alternative is given.
% ---------------------------------------------------------------------------

\usepackage[T1]{fontenc}
\usepackage[utf8]{inputenc}
\usepackage{lmodern}
\usepackage{microtype}
\usepackage[a4paper,margin=2.6cm]{geometry}

\usepackage{amsmath,amssymb,mathtools}
\usepackage{bm}
\usepackage{graphicx}
\usepackage{xcolor}
\usepackage{booktabs}
\usepackage{multirow}
\usepackage{caption}
\usepackage{subcaption}
\usepackage{enumitem}
\usepackage[round]{natbib}
% obeyspaces/spaces let \path print and break at spaces inside evidence notes.
\PassOptionsToPackage{obeyspaces,spaces}{url}
\usepackage[colorlinks=true,linkcolor=blue!55!black,citecolor=blue!55!black,urlcolor=blue!55!black]{hyperref}
\usepackage[capitalize,noabbrev]{cleveref}

% ---------------------------------------------------------------------------
% TODO markers
% ---------------------------------------------------------------------------
\definecolor{TodoRed}{RGB}{190,0,0}
\definecolor{EvidenceGray}{RGB}{110,110,110}

% Inline TODO sentence. Usage at sentence level, never mid-clause.
\newcommand{\todo}[1]{{\color{TodoRed}TODO: #1}}

% Red table cell for pending values.
\newcommand{\tbd}{{\color{TodoRed}TODO}}

% TODO paragraph block.
\newenvironment{TodoBlock}
  {\par\medskip\noindent
   \begin{quote}\color{TodoRed}\small\textbf{TODO:}\enskip\ignorespaces}
  {\end{quote}\medskip}

% Figure slot. A red framed box inside a {figure} float describing the
% figure that has to be produced. Usage \figslot{height}{description}.
\newcommand{\figslot}[2]{%
  \setlength{\fboxsep}{10pt}%
  \setlength{\fboxrule}{1pt}%
  \fcolorbox{TodoRed}{TodoRed!4}{%
    \parbox[c][#1][c]{0.92\linewidth}{\color{TodoRed}\small\centering
      \textbf{TODO: produce this figure.}\\[4pt] #2}}}

% Red caption text for figure slots.
\newcommand{\redcaption}[1]{{\color{TodoRed}TODO: #1}}

% Draft-only evidence notes binding statements to repository artifacts.
\newif\ifevidencenotes
\evidencenotestrue   % switch to \evidencenotesfalse for a clean read
\DeclareRobustCommand{\repopath}[1]{{\small\path{#1}}}
\newcommand{\evidence}[1]{%
  \ifevidencenotes{\color{EvidenceGray}\footnotesize\ [evidence \repopath{#1}]}\fi}

% ---------------------------------------------------------------------------
% Math notation
% ---------------------------------------------------------------------------
\newcommand{\R}{\mathbb{R}}
\newcommand{\E}{\mathbb{E}}
\DeclareMathOperator{\diag}{diag}
\DeclareMathOperator{\rank}{rank}
\DeclareMathOperator{\Sym}{Sym}
\DeclareMathOperator{\softplus}{softplus}
\DeclarePairedDelimiter{\norm}{\lVert}{\rVert}
\newcommand{\T}{^{\mathsf{T}}}

% Frequently used symbols
\newcommand{\field}{F}                 % 2D observation field (teacher)
\newcommand{\student}{S}               % rendered 3D student
\newcommand{\cam}{\pi}                 % projection
\newcommand{\Jq}{J_{\mathrm{Q}}}       % proposal-corrected validation risk
\newcommand{\Ju}{J_{\mathrm{U}}}       % uniform validation risk
\newcommand{\Jpix}{J_{\mathrm{pixel}}}
\newcommand{\Jarea}{J_{\mathrm{area}}}

% Draft typography. Absorb marginal overfulls instead of spilling into the margin.
\emergencystretch=1.5em

% Theorem-like environments
\newtheorem{proposition}{Proposition}
\newtheorem{definition}{Definition}


\title{%
  Compact 2D Gaussian Captures for\\Memory-Efficient 3D Gaussian Splatting\\[6pt]
  {\large with Beam Fusion as a Separately Evaluated Tomographic Initializer}\\[4pt]
  {\normalsize\itshape Working draft --- red text marks what still has to be shown}}

\author{%
  \toshow{Author list t.b.d.}\\
  \toshow{Affiliations t.b.d.}}

\date{Draft of 2026-07-29}

\begin{document}

\maketitle

\begin{abstract}
\noindent
Standard 3D Gaussian Splatting (3DGS) keeps every source image available for the whole
optimization: each iteration renders against randomly selected photographs, so the raw RGB
working set --- not the model --- dominates GPU memory, storage, and bandwidth for dense
multi-camera captures.  We propose to fit every input image \emph{once}, offline, with a compact
per-view 2D Gaussian mixture (a \emph{compact capture}), discard the images, and run all of 3D
reconstruction against the frozen 2D Gaussian fields.  The captures are small, analytic,
differentiable, and can be queried at arbitrary image points, which turns them into streaming
supervision for 3DGS.

The paper makes two deliberately separated arguments.
\textbf{Part~I} carries the systems claim: compact captures contain sufficient information to
supervise reconstruction and refinement of a standard 3DGS representation without ever returning
to source RGB, \toshow{at held-out quality parity with RGB-backed training and with a
$k\times$ reduction in measured peak reconstruction VRAM (headline numbers pending; the
controlled dense-versus-compact resource experiment defined in \cref{sec:vram-protocol} has not
been run yet)}.
\textbf{Part~II} evaluates \emph{Beam Fusion}, an optional image-free initializer that treats the
per-view 2D Gaussians as projections of an unknown 3D radiance density and reconstructs an
initial 3D Gaussian set in closed form by back-projecting each 2D splat into an elongated 3D
beam, gating and fusing beams by covariance intersection, and refining the fused
\emph{carriers} with a fixed-topology compact-only schedule.  Beam Fusion is a research
hypothesis with its own evaluation; by preregistered decision rule, a null or negative Beam
result does not weaken Part~I.

All mechanism-level components are implemented, deterministic on CPU, and covered by an audited
evidence ledger; current quantitative support is single-scene development evidence and is
labelled as such throughout.  \toshow{The confirmatory program --- multi-scene held-out parity,
the controlled VRAM comparison against RGB-backed baselines including Original 3DGS on a real
COLMAP reconstruction, and the matched Beam-versus-no-Beam comparison --- is specified in red
in the corresponding sections and remains to be executed.}
\end{abstract}


\tableofcontents
\bigskip

\section{Introduction}
\label{sec:introduction}

3D Gaussian Splatting \citep{kerbl2023_3dgs} optimizes an explicit set of anisotropic 3D
Gaussians against photographs by differentiable rasterization.  Every iteration samples one
or more source views and compares a full render against the stored photograph, so all
images have to be available during the whole optimization.  Working with dense multi-camera
captures, this becomes the bottleneck.  The source images themselves are the largest
resident data structure of the reconstruction.  A camera dome producing tens to hundreds of
high-resolution views per frame quickly exceeds what a single GPU can hold, and streaming
the images from disk instead trades the memory bottleneck for a storage-bandwidth
bottleneck.  The model being optimized, a few hundred thousand Gaussians, is small by
comparison.

The optimization itself does not consume images.  It consumes colors at projected sample
locations.  Any representation that can answer which color view $v$ sees at image location
$\bm{x}$ can thus supervise the optimization.  We therefore fit every input image once,
offline, with a compact mixture of anisotropic 2D Gaussians, called a compact capture, and
discard the image.  The capture is three orders of magnitude smaller than the raw RGB
working set, see \cref{sec:memory}, analytic, differentiable and can be evaluated at
arbitrary subpixel locations without decoding, resampling or caching full frames.  Our
working hypothesis is that compact 2D Gaussian captures contain sufficient information to
supervise reconstruction and refinement of a standard 3D Gaussian Splatting representation
without returning to source RGB.

The resulting pipeline has a hard, checkable data boundary.  After the per-view fit, no
stage of reconstruction opens an image file, imports an image loader or receives a dense
RGB tensor.  Reconstruction consumes only calibrated cameras plus frozen 2D Gaussian
fields.  We treat this process separation as part of the result and enforce it
structurally through typed inputs, import denials and containment invariants, see
\cref{sec:overview}.\evidence{ara/logic/claims.md C31}

In this paper we therefore describe a supervision method that trains standard 3DGS from
frozen 2D Gaussian fields, together with a compact capture format with exact evaluation
semantics, see \cref{sec:fields,sec:reconstruction}.  The sampled objective is unbiased
and its sparse point renderer matches a dense reference renderer within verified
tolerances.  We define a controlled protocol which measures the memory of dense versus
compact supervision at matched initialization, seeds and stopping rules, see
\cref{sec:memory}.  \todo{Run the protocol and report the memory and quality trade-off.
This is the main missing evidence of the paper, see T-1 to T-6 in \cref{sec:todo-list}.}
Since rendering projects the 3D Gaussian field to per-view 2D Gaussians, the
reconstruction problem also admits a more direct initialization.  Inverting the projection
turns every fitted 2D Gaussian into a constraint on the 3D field.  We describe this
tomographic initialization, called Beam Fusion, in closed form, together with an
observability argument which fixes the minimum number of contributing views at three, see
\cref{sec:tomography}.  The naive inversion fails in identifiable ways.  We analyze these
failures and address them with a carrier refinement consisting of a renderer-aware
covariance repair, a two-phase fixed-topology schedule and a silhouette containment rule,
selected by paired ablations, see \cref{sec:refinement}.  \todo{Replicate the ablations on
fresh scenes with held-out cameras, see T-10.}

We focus on reconstruction from calibrated dense captures.  The reconstruction result is
established with the standard random and SfM initializations alone, so the value of the
tomographic initializer is a separate question with its own evaluation, and a weak
initializer outcome does not touch the reconstruction result.  \Cref{sec:experiments}
evaluates the full pipeline and closes with the ablation study.  Throughout, black text is
mechanism, definition or an evidence-bound result with the binding artifact referenced,
and red text is a TODO.  \Cref{sec:todo-list} collects all TODOs in one table.

\section{Related Work}
\label{sec:related}

\paragraph{3D Gaussian Splatting and its optimization.}
3DGS \citep{kerbl2023_3dgs} represents a scene as anisotropic 3D Gaussians rendered by
EWA-style splatting \citep{zwicker2001_ewa} with front-to-back alpha compositing, optimized
against source photographs with adaptive densification.  A large body of follow-up work
improves the densification and optimization machinery itself: gradient-magnitude corrections
\citep{ye2024_absgs}, pixel-coverage-aware growth \citep{zhang2024_pixelgs}, throughput-focused
re-engineering \citep{mallick2024_taming}, and the reformulation of densification as MCMC
sampling \citep{kheradmand2024_mcmc}.  We deliberately change none of this: both of our
pipeline arms run a standard 3DGS optimizer downstream (\cref{sec:setting}); the manipulated
variables are the \emph{supervision representation} (Part~I) and the \emph{initialization}
(Part~II).

\paragraph{2D Gaussian image representations.}
GaussianImage \citep{zhang2024_gaussianimage} fits images with 2D Gaussians for fast decoding
and compression, establishing that a few thousand anisotropic 2D Gaussians represent a
photograph at high fidelity.  Our Stage~1 (\cref{sec:compact-capture}) adopts this family of
representations as a \emph{capture format for reconstruction}: the fitted field is not a
compression endpoint but the sole supervision signal for 3D optimization.  Note the distinction
from 2D Gaussian \emph{surfel} splatting \citep{huang2024_2dgs}, which changes the 3D scene
primitive; our 2D Gaussians live in image space and describe input photographs.

\paragraph{Compact and compressed 3DGS.}
A separate literature compresses the \emph{output} model: vector quantization and pruning
\citep{lee2024_compact3dgs,fan2024_lightgaussian}, entropy-coded representations
\citep{niedermayr2024_compressed}.  These are orthogonal to our claim, which compresses the
\emph{input side} of the optimization --- the supervision working set held in GPU memory during
training.  The comparison axis of this paper is input representation, not final model size;
output-compression methods compose with our pipeline.

\paragraph{Initialization for 3DGS.}
Standard 3DGS initializes from a COLMAP sparse reconstruction \citep{schoenberger2016_colmap};
random initialization degrades quality or slows convergence, which RAIN-GS
\citep{jung2024_raings} mitigates by schedule design.  Dense learned priors (DUSt3R
\citep{wang2024_dust3r}, MASt3R \citep{leroy2024_mast3r}) power InstantSplat-style pipelines
\citep{fan2024_instantsplat} that initialize from predicted dense geometry.  Beam Fusion
(Part~II) differs in its input: it consumes no images and no learned prior at reconstruction
time, only the fitted 2D Gaussian fields and calibration, and it produces contributor lineage
(which 2D splat in which view supports which 3D Gaussian) as a by-product.  Space-carving
\citep{kutulakos2000_carving} inspires our image-free structural control (compact carve,
\cref{sec:exp2}); covariance intersection \citep{julier1997_ci} supplies the fusion rule and
its consistency guarantee under unknown correlation.

\paragraph{Evaluation conventions.}
We follow community protocol: PSNR, SSIM \citep{wang2004_ssim}, LPIPS \citep{zhang2018_lpips}
on held-out views of standard benchmarks (Mip-NeRF~360 \citep{barron2022_mipnerf360}, Tanks and
Temples \citep{knapitsch2017_tanks}), plus the resource metrics defined in
\cref{sec:vram-protocol}.

\begin{ToShowBlock}[Novelty ledger --- fresh literature pass required before submission]
The field turns over in months; every citation above was written from a stale pass and the
following assumed gaps must each receive a dated verdict with citations before any manuscript
text is frozen (a killed row narrows the claim; discovering it in review kills the paper):
\begin{itemize}[leftmargin=1.4em,itemsep=1pt]
  \item \textbf{N1} --- no published dose--response manipulation of initialization accuracy in
        3DGS (check RAIN-GS successors and 2025/26 initialization-analysis papers);
  \item \textbf{N2} --- no unification of \emph{when} initialization matters across budget,
        view count, and resolution (check sparse-view and budget-constrained 3DGS work);
  \item \textbf{N3} --- no systematic initialization$\times$schedule interaction study (check
        3DGS-MCMC and densification-criteria ablations);
  \item \textbf{N4} --- no published image-free training pipeline (supervising only on
        per-view 2D fits) with quality parity (check compressed/streaming-3DGS and
        feature-space-supervision papers; if this row falls, Part~I's framing shifts from
        ``first'' to the measured $k\times$ numbers);
  \item \textbf{N5} --- dense-prior initializers are evaluated as endpoints, not analyzed for
        \emph{why} they help (InstantSplat lineage; they must appear as baselines here
        regardless).
\end{itemize}
Additionally: verify every bibliography entry against the actual publication record; the
current \repopath{references.bib} was reconstructed from memory and is marked accordingly.
\end{ToShowBlock}

\section{Problem Setting and Pipeline Overview}
\label{sec:setting}

\subsection{Notation}
\label{sec:notation}

A capture session provides $V$ calibrated views.  View $v$ has camera $\cam_v$ (intrinsics
$f_x, f_y$, principal point, world-to-camera rotation $R_v$ and translation), image domain
$\Omega_v = [0, W_v) \times [0, H_v)$, and --- during capture only --- a photograph
$I_v : \Omega_v \to \R^3$.  Stage~1 fits each photograph with a compact 2D Gaussian field
$\field_v$ (\cref{sec:compact-capture}); afterwards the photographs are discarded.  The
reconstruction state is a set of $N$ 3D Gaussians
$\{(\bm{\mu}_i, \Sigma_i, o_i, \bm{c}_i)\}_{i=1}^{N}$ with means
$\bm{\mu}_i \in \R^3$, covariances $\Sigma_i \in \Sym(3)$, opacities $o_i \in (0,1)$, and
SH color coefficients $\bm{c}_i$, rendered exactly as in standard 3DGS.

\begin{definition}[Reconstruction inputs]
\label{def:recon-inputs}
$\mathcal{R} = \{(\cam_v, \field_v)\}_{v=1}^{V}$: calibrated cameras plus frozen per-view 2D
Gaussian observation fields.  No RGB image, alpha mask, or dense tensor is part of
$\mathcal{R}$.
\end{definition}

The data boundary is structural, not aspirational: the reconstruction entry points accept only
$\mathcal{R}$ and configuration; their implementations do not import an image loader or the
dense RGB trainer, do not open image suffixes, and reject dense handovers.  Tests deny these
imports at runtime, so a violation is a test failure, not a code-review judgment
call.\evidence{ara/logic/claims.md C31; src/rtgs/carrier_pipeline.py; tests/test_pipeline.py}

\subsection{The two arms and the RGB control}
\label{sec:arms}

\begin{figure}[t]
  \centering
  \figplaceholder{7.5cm}{%
    \textbf{Figure 1 --- pipeline overview (teaser).}  Three-lane diagram.  Lane 1 (capture,
    grayed out to mark it as upstream of the claim boundary): source RGB $\to$ per-view 2D
    Gaussian fit $\to$ frozen compact captures; RGB shown being discarded (crossed out).
    Lane 2 (Arm 1, no Beam): captures $\to$ Beam-independent 3D initialization (balanced
    compact carve) $\to$ compact-field-supervised 3DGS $\to$ final 3D Gaussians + memory
    receipt.  Lane 3 (Arm 2, Beam): the same captures $\to$ Beam Fusion $\to$ carrier
    refinement $\to$ the same compact optimizer.  A fourth thin lane shows the RGB control
    (standard image-supervised 3DGS from the identical frozen initialization).  Annotate the
    claim boundary (``post-fit reconstruction VRAM'') with a vertical dashed line after the
    fit.  Include per-lane byte sizes for one real scene once measured
    (\cref{sec:vram-protocol}).  Suggested source: draw as TikZ or vector graphics; render
    example thumbnails from \repopath{dataset/} scene frame\_00008 captures.}
  \caption{\redcaption{Pipeline overview and claim boundary.  Figure to be produced; the
    layout is specified in the placeholder.}}
  \label{fig:teaser}
\end{figure}

Both reconstruction arms consume the identical $\mathcal{R}$ and differ only in
initialization; the RGB control isolates the supervision variable.

\textbf{Arm 1 --- direct compact reconstruction (no Beam).}
Balanced compact-carve initialization with a hard 3D count cap
(\cref{sec:exp1-baselines}), then sampled compact-field optimization
(\cref{sec:compact-supervision}) under a frozen no-Beam schedule, then compact-domain
validation and standard \mbox{3DGS} output.  This arm must never import or call Beam Fusion,
Beam contributor lineage, carrier covariance repair, or the Beam-selected carrier schedule.
Its initializer is saved and reused byte-for-byte by the RGB control, so the resource
comparison changes supervision only.  Arm~1 carries \cref{claim:vram}.

\textbf{Arm 2 --- compact reconstruction with Beam Fusion.}
The same captures feed Beam Fusion (\cref{sec:beam-fusion}) and carrier refinement
(\cref{sec:carrier}), followed by the same compact optimizer under matched budgets.  Arm~2
carries only \cref{claim:beam}.

\textbf{RGB control.}
The repository-standard image-supervised 3DGS trainer, loading the exact frozen initialization
produced by Arm~1, with standard densification under a hard cap, on the same device,
resolution, seeds, and stopping rule.  The control may render and score the immutable compact
outputs; the compact processes never see its images.  This control is \emph{not} Original 3DGS:
Original 3DGS requires a real COLMAP sparse reconstruction for initialization, and a random or
Beam-derived initialization must not be relabeled as SfM.
\toshow{Original 3DGS and Original 3DGS with compressed RGB enter as separate baselines once a
real COLMAP model exists for each evaluation scene (obligation O-7, \cref{sec:obligations}).}

\subsection{Predeclared decision rule}
\label{sec:decision-rule}

\begin{itemize}[leftmargin=1.6em,itemsep=1pt]
  \item If Arm 1 reaches the preregistered held-out parity floor and materially lowers peak
        VRAM, the systems claim proceeds regardless of Beam.
  \item If Arm 2 also beats Arm 1 under the matched Beam gate (\cref{sec:exp2-decision}), Beam
        becomes an additional method contribution.
  \item If Arm 2 ties or loses, the direct compact pipeline is retained and Beam is reported as
        a bounded negative or exploratory result; it is not folded into the systems claim.
\end{itemize}

\begin{ToShowBlock}[Rule constants not yet frozen]
The parity floor (dB band on held-out PSNR plus SSIM/LPIPS guards), the materiality threshold
for ``materially lowers peak VRAM'', and the multiplicity policy across scenes and seeds must
be frozen in a preregistration \emph{before} the confirmatory runs, by a reviewer without
outcome access.  The draft blockers listed in
\repopath{docs/THREE_ARM_EXPERIMENT_PROGRAM.md} (exact compact optimizer, topology policy,
stopping rule for Arm 1) must be resolved in the same preregistration.
\end{ToShowBlock}


% ===========================================================================
\paperpart{I}{The Systems Claim: Image-Free Reconstruction at Reduced VRAM}
% ===========================================================================
\section{Compact 2D Gaussian Captures (Stage 1)}
\label{sec:compact-capture}

\subsection{The field model}
\label{sec:field-model}

A compact capture for view $v$ is a mixture of $N_v$ anisotropic 2D Gaussians with an exact,
self-describing evaluation semantics: the capture \emph{is} a frozen renderer, not a loose
point set.  Component $j$ carries a mean $\bm{m}_j \in \Omega_v$, a rotated anisotropic
covariance
\begin{equation}
  C_j \;=\; R(\theta_j)\,\diag\!\big(s_{j,1}^2,\, s_{j,2}^2\big)\,R(\theta_j)\T
  \;+\; \sigma_{f,j}^2\, I_2,
  \label{eq:cov2d}
\end{equation}
with scales $s_j$, rotation $\theta_j$, and an optional isotropic filter variance
$\sigma_{f,j}^2$ (anti-aliasing floor), a non-negative amplitude $a_j \ge 0$, a base color
$\bm{c}_j \in \R^3$ (intentionally unbounded), and optionally a local affine color model
$\bm{c}_j(\bm{x}) = \bm{c}_j + G_j\T(\bm{x} - \bm{m}_j)$ with per-component gradients
$G_j \in \R^{2\times 3}$.  With the unnormalized Gaussian kernel
$w_j(\bm{x}) = a_j \exp\!\big({-\tfrac12 (\bm{x}-\bm{m}_j)\T C_j^{-1} (\bm{x}-\bm{m}_j)}\big)$
truncated at a $3\sigma$ support box, the field renders by normalized blending:
\begin{equation}
  \field_v(\bm{x})
  \;=\;
  \frac{\sum_j w_j(\bm{x})\, \bm{c}_j(\bm{x})}
       {\varepsilon + \sum_j w_j(\bm{x})},
  \qquad \varepsilon = 10^{-8}.
  \label{eq:field}
\end{equation}
Every constant in \cref{eq:cov2d,eq:field} (cutoff, $\varepsilon$, blend mode, fit window,
producer version and configuration digests) is stored inside the capture container, so a
consumer reproduces the fitted image bit-consistently without any out-of-band
convention.\evidence{src/rtgs/core/observation2d.py}

\paragraph{Gauge freedom is an interface hazard.}
The normalized blend makes the pair $(a_j, \bm{c}_j)$ non-identifiable up to
product-preserving rescalings: transformations that leave every rendered color unchanged can
still change downstream algorithms that read raw amplitudes.  We measured exactly this
failure mode: 54 product-preserving Stage-1 re-gaugings left all source renders equivalent
while materially changing coverage and retention of two downstream
lifters.\evidence{ara/logic/claims.md C15}  The paper's rule is therefore: downstream stages
consume \emph{rendered} quantities (field queries, normalized weights) or fix an explicit
gauge; no stage may attach physical meaning to a raw amplitude.  In particular a fitted 2D
amplitude is \emph{not} a 3D opacity (this retires two repair stages in \cref{sec:carrier}).

\subsection{Fitting and the capture container}
\label{sec:fitting}

Stage~1 fits each view independently --- embarrassingly parallel across images --- by
differentiable 2D splatting in the style of GaussianImage \citep{zhang2024_gaussianimage},
with two interchangeable backends in our implementation (a native accumulated-splatting fitter
and a structured fitter with the affine color basis of \cref{eq:field}).  The fit is a
standard L2/robust photometric minimization over all component parameters with a fixed
component budget; masked captures restrict the fit window to the subject.  After fitting, the
field is \emph{frozen}: every later stage treats it as an immutable measurement.

Captures serialize to a small container (one file per view) holding the parameter arrays plus
integrity digests; the current capture format caps at
$168\,\mathrm{kB}$ per view (protocol E11).\evidence{benchmarks/results/20260725_init_value_program_PREREG.md}
For the development scene used throughout (\cref{sec:exp-setup}): native resolution
$5328 \times 4608$, so one raw float32 frame is $294.6$~MB and one 8-bit frame $73.7$~MB,
against $\le 0.168$~MB per capture --- a $440\times$ (8-bit) to $1750\times$ (float32)
reduction of the static input working set before any measurement subtleties.
\toshow{These are format arithmetic, not the claim; the claim-bearing numbers are the measured
process peaks of \cref{sec:vram-protocol}.}

\begin{ToShowBlock}[Stage-1 statistics to report]
A capture-quality table over all evaluation scenes: components per view $N_v$, bytes per view,
fit PSNR/SSIM on the fitted image, fit wall time per view, and backend/configuration.  Plus
the final choice and citation of the Stage-1 fitter configuration (native vs.\ structured
backend, component budget, iteration budget), frozen before the confirmatory runs.  Fit
quality upper-bounds everything downstream (the teacher ceiling), so this table is
load-bearing for interpreting Part~I parity results.
\end{ToShowBlock}

\begin{figure}[t]
  \centering
  \figplaceholder{6cm}{%
    \textbf{Figure 2 --- what a compact capture looks like.}  For two representative views:
    (a) source photograph; (b) rendered compact field $\field_v$ at native resolution; (c)
    signed error map with PSNR inset; (d) ellipse visualization of the fitted 2D Gaussians
    (color = base color, ellipse = $1\sigma$ of $C_j$), optionally with a zoom crop showing
    anisotropic components following image structure.  Data: existing captures under
    \repopath{dataset/}; render (b)/(d) with the repository viewer/preview tooling.}
  \caption{\redcaption{Compact captures: source, field render, error, and component
    structure.  To be produced from the checked-in captures.}}
  \label{fig:capture}
\end{figure}

\subsection{Rate--distortion behaviour}
\label{sec:rate-distortion}

\begin{figure}[t]
  \centering
  \figplaceholder{5cm}{%
    \textbf{Figure 3 --- capture rate--distortion curve.}  Fit PSNR (y) versus bytes per view
    (x, log scale) as the component budget sweeps; one curve per backend.  Add horizontal
    reference lines for JPEG/WebP/AVIF at matched byte budgets to show where the analytic
    representation sits relative to conventional codecs (it does not need to win as a codec
    --- it needs to be \emph{queryable and differentiable} at a competitive rate; make that
    argument in the caption).  Data: sweep Stage-1 fits on 3--5 views of each evaluation
    scene.}
  \caption{\redcaption{Rate--distortion of compact captures versus conventional codecs.
    To be produced.}}
  \label{fig:rd}
\end{figure}

Conventional codecs decode to full frames: supervision from a JPEG still costs a dense RGB
tensor in device memory.  The compact capture inverts the trade: it is competitive on bytes,
\emph{and} its evaluation is a closed-form differentiable function at exactly the sample
locations the optimizer requests, which is what makes the streaming supervision of
\cref{sec:compact-supervision} possible.

\section{Compact-Field Supervision of 3DGS}
\label{sec:compact-supervision}

The optimizer never sees an image; it compares two point-evaluable functions.  The teacher is
the frozen field $\field_v$ of \cref{eq:field}.  The student is the current 3D Gaussian set,
rendered at arbitrary image points by a sparse \emph{point rasterizer}.

\subsection{Point-sampled student rendering}
\label{sec:point-rasterizer}

Given sample locations $\bm{x} \in \Omega_v$, the point rasterizer projects the 3D Gaussians
into view $v$ exactly as the dense renderer does (EWA projection
$P_{iv} = J_{iv} \Sigma_i J_{iv}\T + 0.3\,I_2$, including the $0.3$-pixel dilation), sorts by
view depth, and alpha-composites front to back at each requested point:
\begin{equation}
  \student_v(\bm{x}) \;=\; \sum_{i} T_i(\bm{x})\, \alpha_i(\bm{x})\, \bm{c}_i(\bm{d}_{iv}),
  \qquad
  T_i(\bm{x}) = \prod_{i' < i} \big(1 - \alpha_{i'}(\bm{x})\big),
  \label{eq:student}
\end{equation}
with $\alpha_i(\bm{x}) = o_i \exp(-\tfrac12 q_{iv}(\bm{x}))$, $q_{iv}$ the Mahalanobis form of
$P_{iv}$, and view-dependent SH color $\bm{c}_i(\bm{d}_{iv})$.  The point path matches the
dense CPU reference renderer on color, alpha, depth, global visible order, five 3D
parameter-gradient families, and retained screen-space gradients within sealed tolerances, on
both synthetic grids and a calibrated 4{,}096-point replacement draw; this parity is the
correctness anchor that lets us reason about sampled training as subsampled dense
training.\evidence{ara/logic/claims.md C18;
benchmarks/results/20260716_point_rasterizer_parity_RESULT.json}
\toshow{Boundary: full-image calibrated parity, CUDA/gsplat point parity, and active off-grid
coordinate gradients are not yet established and must be either verified or explicitly scoped
out in the final text.}

\subsection{An unbiased sampled objective}
\label{sec:objective}

Let $p$ be the target risk measure on $\Omega_v$ (uniform over pixels, or a continuous area
measure), and let $q$ be a proposal distribution.  Training draws a \emph{fixed} number $A$ of
attempts $\bm{x}_1, \dots, \bm{x}_A \sim q$ per step and minimizes the importance-corrected
squared error with rejected attempts retained in the denominator:
\begin{equation}
  \widehat{L}_v \;=\;
  \frac{1}{A} \sum_{a=1}^{A}
  \mathbf{1}[\bm{x}_a \text{ accepted}]\;
  \frac{p(\bm{x}_a)}{q(\bm{x}_a)}\;
  \norm{\student_v(\bm{x}_a) - \field_v(\bm{x}_a)}_2^2 .
  \label{eq:sampled-loss}
\end{equation}
Keeping the denominator at the fixed attempt count $A$ (rather than the accepted count) is
what makes \cref{eq:sampled-loss} an unbiased estimator of the target risk under
rejection-mixture proposals; we verified the estimator identity by exact enumeration on finite
pixel grids, including invariance to attempt micro-chunking within
$2\times10^{-12}$.\evidence{ara/logic/claims.md C19}

Two evaluation risks are computed \emph{exactly} and are never replaced by the training loss:
$\Jpix$, the mean over all pixel centers, and $\Jarea$, a fixed four-offset quadrature per
pixel.  In experiment tables we report the validation risks $\Ju$ (uniform sampling) and
$\Jq$ (proposal-corrected), both against the frozen fields of held-out-for-validation
views.\evidence{benchmarks/compact_only_carrier_stage_ablation.py}

\paragraph{Proposal choice is settled empirically.}
Teacher-shaped Gaussian proposals are the obvious ``importance'' candidate; in a sealed
three-seed fixed-topology comparison they produced no material win --- the discrete-pixel
Gaussian mixture lost direction in all three seeds (geometric-mean AUC ratio $1.0246$), and
the continuous-area Gaussian mixture was directionally favorable in all seeds
($0.9911$) but missed the preregistered $0.95$ materiality floor.  The pipeline therefore
defaults to uniform proposals; this negative result is part of the method's evidence
base.\evidence{ara/logic/claims.md C20}

\subsection{Streaming supervision and the memory shape}
\label{sec:streaming}

Because both sides of \cref{eq:sampled-loss} are point queries, a training step touches:
the compact fields of the sampled views (kilobytes each), the sampled coordinate batch, the
projected student parameters, and the model/optimizer state.  No dense render target, no
image cache, and no all-view RGB tensor exist anywhere in the loop --- the fields are queried
on demand and the hard boundary of \cref{sec:arms} forbids materializing dense frames.  This
is the mechanism behind \cref{claim:vram}; what it buys is quantified (only) by the protocol
of \cref{sec:vram-protocol}.

\begin{ToShowBlock}[Supervision-side obligations]
(i) Freeze the Arm-1 optimizer configuration (sample budget per step, view schedule,
learning rates, topology policy, stopping rule) in the preregistration --- these are draft
blockers, deliberately not chosen in this text.
(ii) Show a sampled-versus-dense equivalence curve: quality of compact-supervised training as
a function of samples per step, against dense-RGB training at matched update counts
(\cref{fig:sample-budget}).
(iii) Establish or scope out CUDA point-rasterizer parity (the CPU reference anchors
correctness; throughput claims need the CUDA path).
\end{ToShowBlock}

\begin{figure}[t]
  \centering
  \figplaceholder{5cm}{%
    \textbf{Figure 4 --- sample-budget sweep.}  Held-out PSNR (y) versus samples per step
    (x, log scale) for compact-field supervision, with the dense-RGB control as a horizontal
    band and matched update counts.  One panel per scene; seeds as thin lines, median bold.
    Purpose: locate the budget where sampled supervision saturates, justifying the chosen
    default.  Data: Arm-1 training sweeps once the Arm-1 schedule is frozen.}
  \caption{\redcaption{Quality versus sample budget for compact supervision.  To be
    produced.}}
  \label{fig:sample-budget}
\end{figure}

\section{The VRAM Claim and Its Measurement Protocol}
\label{sec:vram-protocol}

\subsection{Claim boundary}
\label{sec:claim-boundary}

The honest name of the measured quantity is \emph{post-fit reconstruction VRAM}.  The
reconstruction boundary starts at integrity-checked reconstruction inputs
(\cref{def:recon-inputs}) and ends when the final 3D Gaussian file is saved.  RGB capture and
the RGB$\to$2D fit are upstream of the boundary: they run once, offline, per view,
embarrassingly parallel, and on hardware unconstrained by the reconstruction GPU.  We disclose
the upstream fit's storage, time, and VRAM alongside the claim
(\toshow{table pending, obligation O-2}), but the claim itself concerns the resource profile
of reconstruction --- the stage that today forces dataset-scale GPU memory.

\subsection{Memory model}
\label{sec:memory-model}

Let $b$ be bytes per channel, $V$ the view count, $H \times W$ the resolution, $N$ the 3D
Gaussian count.  To first order:
\begin{align}
  M_{\text{dense}} &\;\approx\; \underbrace{V \, H \, W \, 3\, b}_{\text{resident RGB working set}}
    \;+\; M_{\text{model+opt}}(N) \;+\; M_{\text{render}}(H, W),
  \label{eq:mem-dense}\\[2pt]
  M_{\text{compact}} &\;\approx\; \underbrace{\textstyle\sum_{v} N_v \, \beta}_{\text{all captures}}
    \;+\; M_{\text{model+opt}}(N) \;+\; M_{\text{query}}(A),
  \label{eq:mem-compact}
\end{align}
with $\beta \approx 40$ bytes per 2D component and $M_{\text{query}}$ the per-step sampled
batch (\cref{sec:streaming}), which is independent of $V$, $H$, $W$.  The dense working set
grows linearly in view count and quadratically in resolution; the compact working set is
capped by the capture budget ($\le 168$~kB/view here) and is dwarfed by
$M_{\text{model+opt}}$.  For the development scene ($V{=}26$, $5328 \times 4608$):
$7.7$~GB (float32) or $1.9$~GB (uint8) of resident RGB versus $4.4$~MB of captures.
Streaming dense images from disk instead of holding them trades
$M_{\text{dense}}$ for IO bandwidth and latency; the protocol therefore measures IO alongside
memory.

\Cref{eq:mem-dense,eq:mem-compact} are an accounting sketch, not evidence: real allocators
fragment, renderers hold workspaces, and CUDA context overhead is shared.  The claim is
decided exclusively by the measured protocol below.

\subsection{Measurement protocol}
\label{sec:protocol}

Each measured cell runs in a \emph{fresh process} on an otherwise idle, named GPU and
records:
\begin{itemize}[leftmargin=1.6em,itemsep=1pt]
  \item NVML process-used-memory peak, sampled at a frozen rate (the headline number;
        CUDA-allocator counters miss driver/backend allocations and cannot be the sole VRAM
        figure);
  \item PyTorch allocated and reserved peaks (for decomposition, reported separately);
  \item process RSS (host RAM), bytes read/written, wall time and per-stage time, final
        Gaussian count;
  \item software, driver, and GPU identity, plus the background-memory baseline before
        process start.
\end{itemize}
Rules: one warm-up run, then $\ge 3$ measured repeats ($\ge 5$ for throughput) per
seed$\times$scene cell; identical device, resolution, initialization (byte-for-byte,
\cref{sec:arms}), seeds, and stopping rule across arms; A/B-interleaved execution order to
cancel thermal drift; every raw repeat retained --- headline medians alone are insufficient;
the compact working-set size reported separately from model/optimizer memory.  Headline
metrics:
\begin{enumerate}[leftmargin=1.6em,itemsep=1pt]
  \item peak reconstruction VRAM, compact versus RGB-backed, at matched primitive and view
        counts (the reduction factor $k$);
  \item the maximum number of simultaneously-held views at fixed VRAM budget (``$n$ views on
        one GPU versus $n/k$ RGB-backed'');
  \item bytes on disk per view; iterations per second at matched primitive count.
\end{enumerate}
The predeclared claim form is: \emph{at fixed VRAM budget and fixed primitive count, the
compact path holds $k\times$ more views than the RGB path, measured as above, at held-out
quality inside the preregistered parity band.}  Quality parity is the load-bearing word:
``slightly worse but smaller'' is a different and much weaker paper, and the decision rule of
\cref{sec:decision-rule} treats it as such.

\begin{ToShowBlock}[The entire quantitative content of this section is pending]
No cell of this protocol has been run.  Required: (i) the NVML sampler and fresh-process
resource wrapper shared by compact and RGB arms; (ii) the frozen parity floor and materiality
rule (\cref{sec:decision-rule}); (iii) at least one outcome-unseen calibrated scene ---
the development frames used so far are outcome-exposed and cannot become confirmatory by
renaming; (iv) standard benchmark scenes (Mip-NeRF~360, Tanks and Temples) for external
validity; (v) the $k\times$ and views-at-fixed-VRAM measurements themselves
(\cref{fig:vram-views,tab:main-results}).  Everything stated above it is protocol design,
implemented mechanism, or format arithmetic.
\end{ToShowBlock}

\begin{figure}[t]
  \centering
  \figplaceholder{5.5cm}{%
    \textbf{Figure 5 --- the headline figure.}  Peak reconstruction VRAM (y) versus number of
    views (x) at fixed primitive count: one line for RGB-backed training (grows linearly,
    hits the GPU capacity line), one for compact supervision (near-flat).  Mark the GPU
    capacity as a horizontal line; annotate the crossing point (``RGB path exceeds a 24~GB
    GPU at $V{=}\dots$'') and the $k\times$ gap at the largest common $V$.  Optionally a
    second panel: iterations/second versus views.  Data: the protocol of
    \cref{sec:protocol} on one dome scene with view-count sweeps; NVML peaks with repeat
    spread as error bars.}
  \caption{\redcaption{Peak VRAM versus view count, dense versus compact.  The claim-carrying
    figure; to be produced from the controlled resource experiment.}}
  \label{fig:vram-views}
\end{figure}

\section{Experiments: Part I}
\label{sec:exp1}

\subsection{Setup}
\label{sec:exp-setup}

\paragraph{Data.}
Development used one calibrated multi-camera capture (a masked studio subject,
``stage-with-fabric'' frame~00008): 26 cameras at $5328 \times 4608$, fitted compact captures
for all views, split 19 fitting / 7 compact-validation views.  This scene is
\emph{outcome-exposed} --- it drove every design decision --- and is therefore reported as
development evidence only.
\toshow{Confirmatory data: at least one fresh calibrated dome scene with untouched test
cameras, plus Mip-NeRF~360 and Tanks-and-Temples scenes with the Stage-1 fit run at benchmark
resolution.  Frozen train/validation/test protocol per scene before any confirmatory run;
test cameras touched exactly once for the final tables.}

\paragraph{Metrics.}
Held-out PSNR, SSIM, LPIPS (every reported cell, not PSNR alone), alpha IoU and exterior
leakage for masked subjects; the resource metrics of \cref{sec:protocol}; compact-domain
$\Ju$/$\Jq$ for diagnostics.  Geometry error against reference depth/mesh where a reference
exists --- PSNR is blind to floaters, and floaters are the failure mode initialization is
supposed to prevent.

\subsection{Baselines}
\label{sec:exp1-baselines}

\begin{table}[t]
  \centering\footnotesize
  \setlength{\tabcolsep}{4pt}
  \begin{tabular}{@{}lp{3.1cm}p{3.6cm}p{3.6cm}@{}}
    \toprule
    Arm & Supervision & Initialization & Status \\
    \midrule
    Original 3DGS            & RGB & COLMAP sparse
      & \toshow{needs real COLMAP model} \\
    Orig.\ 3DGS, compr.\ RGB & matched-byte JPEG/AVIF & COLMAP sparse
      & \toshow{pending} \\
    RGB control              & RGB & frozen Arm-1 init (byte-identical)
      & \toshow{pending} \\
    Direct compact (Arm 1)   & compact fields & balanced compact carve
      & \toshow{pending (schedule to freeze)} \\
    Beam + compact (Arm 2)   & compact fields & Beam Fusion + carriers
      & \toshow{pending (Part II gate)} \\
    \bottomrule
  \end{tabular}
  \caption{Arms of the Part-I comparison.  The RGB control isolates the supervision variable
    (identical initialization); the Original-3DGS rows anchor against the community-standard
    pipeline.  \toshow{All quantitative cells pending; a random-initialization RGB control and
    a fixed-topology RGB control are added to separate supervision value from capacity
    growth.}}
  \label{tab:arms}
\end{table}

Arm 1's initializer is the \emph{balanced compact carve}: one bounded ray per fitted 2D
component (source components propose rays and 2D covariances but never supervise), scored by
coverage-weighted queries of every other frozen field along the ray's depth interval, with
per-view balanced top-$K$ selection under a hard global count cap.  Related development
evidence: unrestricted dense placement beat balanced top-$K$ initialization by
$1.97$~dB but violated the count budget by $13.5\times$, and a confidence-gated ``easy-only''
variant stayed $2.17$~dB below dense placement at the matched schedule --- so the balanced
capped variant is the frozen default and count honesty is part of the
protocol.\evidence{ara/logic/claims.md C24}

\subsection{Main results}
\label{sec:exp1-results}

\begin{table}[t]
  \centering\footnotesize
  \setlength{\tabcolsep}{4pt}
  \begin{tabular}{@{}lcccccccc@{}}
    \toprule
    & \multicolumn{3}{c}{Held-out quality} & \multicolumn{3}{c}{Resources (peak)} & & \\
    \cmidrule(lr){2-4}\cmidrule(lr){5-7}
    Arm & PSNR$\uparrow$ & SSIM$\uparrow$ & LPIPS$\downarrow$
        & VRAM & host RAM & disk & IO & time \\
    \midrule
    Original 3DGS            & \tbd & \tbd & \tbd & \tbd & \tbd & \tbd & \tbd & \tbd \\
    Orig.\ 3DGS, compr.\ RGB & \tbd & \tbd & \tbd & \tbd & \tbd & \tbd & \tbd & \tbd \\
    RGB control (same init)  & \tbd & \tbd & \tbd & \tbd & \tbd & \tbd & \tbd & \tbd \\
    Direct compact (Arm 1)   & \tbd & \tbd & \tbd & \tbd & \tbd & \tbd & \tbd & \tbd \\
    Beam + compact (Arm 2)   & \tbd & \tbd & \tbd & \tbd & \tbd & \tbd & \tbd & \tbd \\
    \bottomrule
  \end{tabular}
  \caption{\redcaption{Main Part-I table: held-out quality and measured resource peaks per
    arm (median over $\ge 3$ fresh-process repeats, spread in the appendix; per-scene
    sub-tables).  Entirely pending the controlled experiment.}}
  \label{tab:main-results}
\end{table}

\begin{figure}[t]
  \centering
  \figplaceholder{6cm}{%
    \textbf{Figure 6 --- qualitative parity grid.}  For 2--3 held-out test cameras per scene:
    RGB-control render, Arm-1 render, difference maps, and crops of the hardest regions
    (hair/fabric edges, specular highlights).  Purpose: make the parity (or its failure)
    inspectable; reviewers will look for exactly these crops.  Data: final models of
    \cref{tab:main-results}.}
  \caption{\redcaption{Qualitative held-out comparison, RGB-backed versus compact-supervised.
    To be produced.}}
  \label{fig:parity-grid}
\end{figure}

\subsection{Existing development evidence (what is already known)}
\label{sec:exp1-dev-evidence}

No cell of \cref{tab:main-results} exists yet; the following adjacent results bound
expectations and are part of the honest record.
A full-scale RGB-consuming development run on the exposed scene (all 26 views fitted, no
held-out cameras) reached $28.76$~dB native foreground PSNR / $0.974$ SSIM / $0.0173$ LPIPS
at 100k Gaussians after 128k updates, confirming that the capture/optimization machinery
scales to native resolution --- but it authorizes no compact-only, held-out, or resource
statement (every optimized phase consumed RGB).\evidence{ara/logic/claims.md C29}
The compact-only carrier pipeline of Part~II reached compact-validation risks
$\Jq = 0.00422$, $\Ju = 0.00786$ on the same scene under a strict no-image
boundary,\evidence{ara/logic/claims.md C30, C31} establishing that the compact objective
optimizes stably end-to-end.  The mechanism results of \cref{sec:compact-supervision}
(rasterizer parity, estimator identity, proposal null result) are seed-replicated CPU
evidence.  None of this substitutes for \cref{tab:main-results}.

\begin{ToShowBlock}[Part-I experiment obligations]
(i) Build the shared NVML/fresh-process resource harness (obligation O-1).
(ii) Freeze Arm-1 schedule, parity floor, materiality rule, multiplicity policy
(O-3).
(iii) Run \cref{tab:main-results} on the fresh dome scene + benchmark scenes with $\ge 3$
paired seeds (O-4/O-5).
(iv) Produce \cref{fig:vram-views} (view-count sweep) and \cref{fig:parity-grid} (O-6).
(v) Acquire real COLMAP reconstructions for the Original-3DGS rows (O-7).
(vi) Report upstream Stage-1 fit cost table (O-2).
\end{ToShowBlock}


% ===========================================================================
\paperpart{II}{Beam Fusion: Tomographic Initialization from Compact Captures}
% ===========================================================================
\section{Beam Fusion}
\label{sec:beam-fusion}

Part II asks whether the compact captures can also produce a \emph{better 3D initialization}
than the structural no-Beam baseline.  The key observation: each per-view 2D Gaussian field is
a \emph{projection} of the same unknown 3D radiance density.  3D initialization can therefore
be posed as reconstruction-from-projections --- tomography --- instead of discrete
correspondence search.  Gaussians make this tractable in closed form: the back-projection of a
2D Gaussian is an analytic elongated 3D Gaussian, and products/fusions of Gaussians remain
Gaussian.  Beam Fusion is deterministic, RGB-free, correspondence-free (association emerges
from density overlap, not discrete matching), and CPU-first.

\subsection{Back-projection: from 2D splats to 3D beams}
\label{sec:beams}

Fix view $v$ with camera center $\bm{o}_v$ and a fitted component $(\bm{m}_j, C_j, a_j,
\bm{c}_j)$.  Its center pixel defines a world ray direction $\bm{d}_{vj}$.  At a hypothesised
camera depth $t$, the projective Jacobian of the pinhole map at the point
$\bm{p} = t\,\bm{r}_{vj}$ (camera frame) is
\begin{equation}
  J(\bm{p}) \;=\;
  \begin{pmatrix}
    f_x / z & 0 & -f_x\, x / z^2\\
    0 & f_y / z & -f_y\, y / z^2
  \end{pmatrix},
  \qquad \bm{p} = (x, y, z)\T .
  \label{eq:jacobian}
\end{equation}
Let $B = (\bm{b}_1, \bm{b}_2, \bm{r})$ be an orthonormal camera-frame basis whose third column
is the ray direction.  The fitted 2D covariance is lifted to the ray-orthogonal tangent plane
through the inverse of the projected tangent Jacobian,
\begin{equation}
  \Sigma_{\perp}
  \;=\;
  \big(J B_{:,1:2}\big)^{-1}\, C_j\, \big(J B_{:,1:2}\big)^{-\mathsf{T}},
  \label{eq:tangent-lift}
\end{equation}
which is exact for the center ray: it answers ``which tangent-plane covariance would have
projected to the observed $C_j$?''.  Along the ray the component's depth is unknown within
the working depth range $[t_{\mathrm{lo}}, t_{\mathrm{hi}}]$ obtained by intersecting the ray
with the scene bounding box (scaled camera-cloud extent, clamped to a near plane); the beam
assigns it a long variance $L^2$ with half-length
$L = \tfrac12 (t_{\mathrm{hi}} - t_{\mathrm{lo}})$.  The \emph{beam} is the 3D Gaussian with
precision
\begin{equation}
  \Lambda_{vj}(t)
  \;=\;
  R_v\T \, B \begin{pmatrix} \Sigma_{\perp}^{-1} & 0 \\ 0 & 1/L^2 \end{pmatrix} B\T R_v
  \label{eq:beam-precision}
\end{equation}
in world coordinates: certain transversally exactly as the image measured, agnostic along the
ray.  No voxel grid is ever built.

\subsection{Pair seeding in closed form}
\label{sec:seeding}

For components $i$ in view $u$ and $j$ in view $v$, the closest points between their center
rays have closed-form depths $(s, t)$ (the standard two-ray least-squares solution; parallel
rays are rejected).  A candidate 3D location is the midpoint of the closest-point segment.
Candidates are gated by three tests, all analytic:
\begin{equation}
  \underbrace{\frac{\norm{\bm{p}_u(s) - \bm{p}_v(t)}^2}
       {\sigma_u^2(s) + \sigma_v^2(t)} \le 3^2}_{\text{transverse gate}},
  \qquad
  \underbrace{\norm{\bm{c}_i - \bm{c}_j} \le 0.35}_{\text{color gate}},
  \qquad
  \underbrace{s \in [t^u_{\mathrm{lo}}, t^u_{\mathrm{hi}}],\;
              t \in [t^v_{\mathrm{lo}}, t^v_{\mathrm{hi}}]}_{\text{depth validity}},
  \label{eq:gates}
\end{equation}
where $\sigma_u(s)$ is component $i$'s mean pixel footprint scaled to world units at depth $s$
--- i.e.\ the ray separation is measured \emph{against both beams' transverse footprints}, so
large splats tolerate larger separations.  Each surviving seed receives the first-order
Gaussian product weight
\begin{equation}
  w \;=\; \tfrac12 (a_i + a_j)\,
          \exp\!\big({-\tfrac12\, g}\big)\,
          \exp\!\big({-\norm{\bm{c}_i - \bm{c}_j}^2 / (2\sigma_c^2)}\big),
  \qquad \sigma_c = 0.25,
  \label{eq:seed-weight}
\end{equation}
with $g$ the transverse gate value: negligible cross terms are never materialized.  View
pairs are processed in increasing camera-distance order; the production path streams all
pairs through a bounded 3D voxel table that retains only the strongest seed per voxel
(deterministic tie-breaking by lexicographic rank), so seed storage is capped independent of
scene size.

\subsection{Fusion by covariance intersection}
\label{sec:ci-fusion}

Contributing beams of one physical splat are \emph{correlated} observations --- they look at
the same surface through overlapping information.  The naive Gaussian product
$\Lambda = \sum_k \Lambda_k$ is therefore rejected by design: it double-counts every axis
observed by more than one view (two orthogonal views of an isotropic blob would fuse to half
its variance).  Beam Fusion uses equal-weight covariance intersection
\citep{julier1997_ci}:
\begin{equation}
  \Lambda \;=\; \frac{1}{K} \sum_{k=1}^{K} \Lambda_k,
  \qquad
  \bm{m} \;=\; \Lambda^{-1} \Big( \frac{1}{K} \sum_{k=1}^{K} \Lambda_k\, \bm{m}_k \Big),
  \label{eq:ci}
\end{equation}
where $\bm{m}_k$ is beam $k$'s world-space mean at its implied depth.  The common $1/K$
leaves the fused mean identical to the product rule but keeps the covariance conservative:
CI is exact on directions all views share and never overconfident elsewhere, which is the
standard consistency guarantee under unknown correlation.  Consequently Beam Fusion's main
failure risk is \emph{false contributor association}, not an algebraic missing factor ---
which is why contributor counts, gates, and lineage are first-class outputs
(\cref{sec:reduction}).

\subsection{Why at least three views: an observability result}
\label{sec:observability}

\begin{proposition}[Two views leave one covariance coordinate unobservable]
\label{prop:rank5}
Fix a 3D Gaussian with triangulated mean and EWA covariance projection
$C_v = A_v \Sigma A_v\T$ into views $v = 1, 2$ with world ray directions
$\bm{d}_1, \bm{d}_2$ through the mean.  The stacked linear map
$\Sigma \mapsto (C_1, C_2)$ on $\Sym(3) \cong \R^6$ has rank five for generic camera pairs,
and the symmetric matrix $\bm{d}_1 \bm{d}_2\T + \bm{d}_2 \bm{d}_1\T$ spans a shared null
direction.  Three generic views raise the rank to six.
\end{proposition}

\emph{Proof sketch.}  Each projection Jacobian annihilates its own viewing ray,
$A_v \bm{d}_v = 0$.  For the null candidate $\Delta = \bm{d}_1 \bm{d}_2\T + \bm{d}_2
\bm{d}_1\T$: $A_1 \Delta A_1\T = (A_1\bm{d}_1)(A_1\bm{d}_2)\T + (A_1\bm{d}_2)(A_1\bm{d}_1)\T
= 0$, and symmetrically for view 2, so $\Delta$ is invisible to both views; a dimension count
over $\Sym(3)$ gives generic rank five with exactly this shared mode, and a third generic ray
breaks it.  The full statement is verified symbolically and numerically in the repository
test suite.\evidence{ara/logic/claims.md C23;
tests/test_inverse_projection_fiber.py}\hfill$\square$

Two-view fusions would therefore carry an arbitrary along-baseline covariance component ---
not merely a noisy one.  Beam Fusion requires $K \ge 3$ contributing views per output
component (\texttt{min\_views}=3), turning an algebraic blind spot into a hard admission
rule.

\subsection{Greedy fold-in, reduction, and lineage}
\label{sec:reduction}

Each surviving seed projects its provisional 3D midpoint into every remaining view and
absorbs at most one nearest fitted component per view, subject to a $3\sigma$ pixel gate
(relative to the candidate component's own footprint), the color gate of \cref{eq:gates},
and a positive-depth test.  Components with fewer than three contributors are discarded
(\cref{prop:rank5}).  Each retained component is CI-fused (\cref{eq:ci}) over all its
contributors, weighted $w \cdot K$, and reduced in two exact steps: (i) contributor-signature
dedupe (identical view/component sets collapse to the strongest instance), and (ii) per-voxel
non-maximum suppression by weight --- \emph{selection}, not moment matching, so fused
covariances survive unblurred.  Covariance eigenvalues are clamped to
$[\sigma_{\min}^2, \sigma_{\max}^2]$ world bounds; colors initialize as amplitude-weighted
contributor means; opacity initializes at a constant $0.1$ (\cref{sec:carrier} explains why
no opacity can be inferred from the captures).  The output is a standard 3D Gaussian set
\emph{plus complete contributor lineage} (a CSR map from every 3D component to its
supporting 2D splats and implied depths) and per-view unmatched-splat lists --- the
correspondence by-product that downstream refinement, diagnostics, and future densification
consume.  Diagnostics recorded per run: contributor-count histogram, evaluated ray pairs,
gate survival counts, seed-voxel occupancy, and stage timings.\evidence{src/rtgs/lift/beam\_fusion.py}

\begin{figure}[t]
  \centering
  \figplaceholder{6.5cm}{%
    \textbf{Figure 7 --- Beam Fusion geometry.}  Two-panel schematic.  (a) One 2D Gaussian in
    two views back-projected as two elongated beams (transverse footprint from
    \cref{eq:tangent-lift}, long axis $L$ from the depth interval); their closest-ray
    midpoint, the transverse gate as shaded tolerance tubes, and the CI-fused 3D Gaussian at
    the intersection --- visually contrast CI (conservative ellipsoid) with the naive
    product (overconfident small ellipsoid).  (b) Fold-in: the provisional point projecting
    into a third view, the $3\sigma$ pixel gate around the nearest 2D component, and the
    resulting 3-contributor lineage.  Suggested: TikZ 3D schematic; exact ellipse shapes can
    be computed from a toy two-camera configuration.}
  \caption{\redcaption{Beam back-projection, gating, and covariance-intersection fusion.
    Schematic to be drawn.}}
  \label{fig:beam-geometry}
\end{figure}

\begin{figure}[t]
  \centering
  \figplaceholder{6cm}{%
    \textbf{Figure 8 --- initialization comparison (visual).}  Renders (front + side/top
    view) of the initial 3D Gaussian sets on one scene: Beam Fusion, balanced compact carve,
    calibration-bounded random, and (if available) COLMAP sparse init; same count budget.
    Side/top views expose along-ray elongation and floaters that front views hide.  Include
    contributor-count coloring for the Beam panel (3, 4, $\ge$5 views).  Data: existing
    development artifacts under \repopath{runs/} or a rerun of the three initializers on the
    development scene; render with \repopath{rtgs view} / preview tooling.}
  \caption{\redcaption{Initial 3D Gaussians per initializer, front and side views.  To be
    produced from checked-in initializers.}}
  \label{fig:init-visual}
\end{figure}

\section{Carrier Refinement}
\label{sec:carrier}

Beam Fusion is deliberately \emph{not} a drop-in initializer.  Its means are accurate early,
but covariance, opacity, and appearance start from geometry-only reasoning; dropped directly
into a standard 3DGS optimizer with densification, the optimizer tends to \emph{replace}
the initialization rather than refine it --- and the experiment then measures densification's
ability to recover, not the initializer.  Beam outputs are therefore treated as
\emph{carriers}: they pass through a short, fully compact-only maturation before any
comparison.  The accepted carrier sequence was selected by three preregistered, paired-root,
independently audited ablations\evidence{docs/ADR-002-carrier-refinement.md;
ara/logic/claims.md C30, C31} and consists of five stages:
\begin{enumerate}[leftmargin=1.6em,itemsep=1pt]
  \item fit-only Beam Fusion from reconstruction inputs ($K \ge 3$, no free primitive
        birth);
  \item renderer-aware covariance repair (\cref{sec:cov-repair});
  \item 380 compact fixed-topology SH0 updates with means, rotations, scales, opacity, and
        SH0 all trainable;
  \item strict projected-center pruning against every fitting view (\cref{sec:containment});
  \item 380 further compact SH0 updates with means frozen, so the pruning invariant cannot
        be undone.
\end{enumerate}
No clone, split, insertion, densification, opacity reset, higher-SH phase, or dense handover
exists on this path, and its entry point accepts no image-bearing argument.

\subsection{Renderer-aware covariance repair}
\label{sec:cov-repair}

The point renderer projects a 3D covariance as $P_{iv} = J_{iv} \Sigma_i J_{iv}\T + 0.3\,
I_2$ (\cref{sec:point-rasterizer}); any repair target that omits the $0.3$-pixel EWA dilation
is inconsistent with what the renderer will actually display.  For carrier $i$ with
contributor set $\mathcal{V}_i$ (from lineage) and observed 2D covariances $C_{iv}$, define
the whitened projection
\begin{equation}
  M_{iv} \;=\; C_{iv}^{-1/2}\,\big(J_{iv}\, \Sigma_i\, J_{iv}\T + 0.3\, I_2\big)\,
               C_{iv}^{-1/2},
  \label{eq:whitened}
\end{equation}
and minimize, over $\Sigma_i$ only (means, opacity, appearance, lineage, and count stay
bit-frozen), the symmetric generalized-log residual
\begin{equation}
  r_i \;=\; \sqrt{\;\frac{1}{|\mathcal{V}_i|}\sum_{v \in \mathcal{V}_i}
        \operatorname{mean}_k \log^2 \lambda_k(M_{iv})\;},
  \label{eq:cov-residual}
\end{equation}
with a Huber data term, a prior toward the Beam covariance, and explicit SPD/aspect bounds.
The $\log$-eigenvalue form is the load-bearing choice: it penalizes a projection twice as
large and half as large equally.  The previous repair minimized a one-sided whitened
Frobenius residual without the dilation term; that objective penalizes expansion more than
collapse, is shrink-biased, and was empirically harmful --- it worsened factorial-marginal
validation $\Jq$ by $51.6\%$.  The corrected objective of
\cref{eq:whitened,eq:cov-residual} reduced validation $\Jq$ by $63.2\%$ versus directly
optimizing raw Beam output, with 3/3 paired-root wins.\evidence{ara/logic/claims.md C30}
Because some fitted targets are sharper than the renderer's $0.3$-pixel floor, zero residual
is not always attainable; the repair aims at consistency, not perfection.

\paragraph{What is deliberately \emph{not} repaired.}
Opacity: a fitted 2D amplitude is a normalized-blend gauge quantity
(\cref{sec:field-model}), not an identifiable per-carrier 3D alpha under compositing ---
legacy opacity repair worsened $\Jq$ by $9.8\%$ and is retired.  Appearance: cross-view
amplitude-weighted color averaging has the same gauge problem and was immaterial ($0.12\%$,
below the $5\%$ necessity floor); SH0 colors are instead co-adapted by the compact phases.
Topology: even an optical-density-preserving clone (copying opacity changes coincident alpha
from $a$ to $1-(1-a)^2$; the preserving variant uses $a' = 1 - \sqrt{1-a}$) failed its stage
gate; the carrier path grows nothing.\evidence{ara/logic/claims.md C30;
docs/ADR-002-carrier-refinement.md}

\subsection{Two compact phases; why all parameter families train together}
\label{sec:phases}

Both phases minimize the sampled compact objective of \cref{eq:sampled-loss} at fixed
topology and SH degree 0.  Phase 1 trains all families jointly: the alpha/color gauge freedom
of the compositor means opacity and support cannot be identified independently --- scale,
orientation, opacity, and color must co-adapt.  Means-only optimization was strongly
rejected ($4.18\times$ the all-family phase-1 $\Jq$); a second all-family phase reduced
$\Jq$ to $0.717\times$ its stopping value with 3/3 wins; freezing \emph{only the means} in
phase 2 costs $1.47\%$ $\Jq$ and passes the frozen $2\%$ non-inferiority gate --- which is
exactly what makes the containment invariant of \cref{sec:containment} permanent.
Incremental SH3 improved $\Jq$ by only $2.3$--$3.0\%$, below the $5\%$ materiality floor, so
the carrier stays SH0.\evidence{ara/logic/claims.md C30}

\subsection{Projected-center containment}
\label{sec:containment}

Raw masks are forbidden inside the compact arms; their legal surrogate is the union of fitted
2D Gaussian support ellipses.  After phase 1, carrier $i$ is retained only if, for
\emph{every} fitting view $v$:
\begin{equation}
  z_{iv} > z_{\text{near}}
  \qquad\text{and}\qquad
  \min_j \big(\cam_v(\bm{\mu}_i) - \bm{m}_{vj}\big)\T C_{vj}^{-1}
         \big(\cam_v(\bm{\mu}_i) - \bm{m}_{vj}\big) \;\le\; 3^2,
  \label{eq:containment}
\end{equation}
i.e.\ the projected center must fall inside some positive-amplitude fitted component's
$3\sigma$ ellipse in every view that observes the scene.  The threshold is fixed and cannot
be weakened through configuration; phase 2 freezes means and rechecks
\cref{eq:containment}, so the invariant survives to the output.  On the development scene the
prune removes $5.3$--$5.4\%$ of carriers, retains $4{,}729$--$4{,}734$, ends with exactly
zero violations, and costs $2.69\%$ $\Jq$ against the unpruned
anchor.\evidence{ara/logic/claims.md C30, C31}

Honesty about what \cref{eq:containment} is: a \emph{visual-hull center invariant}.  It
removes carriers whose centers leave the fitted silhouette cone; it cannot detect a floater
\emph{inside} the multi-view hull, and it does not prove surface occupancy.  Whole-Gaussian
mask containment is impossible in principle (Gaussian support is infinite), and a
finite-footprint rule would erode legitimate silhouette-boundary primitives.  A
count-preserving soft variant is predeclared for the ablation grid: the barrier
\begin{equation}
  L_{\text{support}} = \sum_{i,v}
    \softplus\!\Big(\tfrac{\min_j q_{ivj} - 9}{\tau}\Big)^{2}
    + \lambda_z\, \softplus\!\Big(\tfrac{z_{\text{near}} - z_{iv}}{\tau_z}\Big)^{2},
  \label{eq:soft-support}
\end{equation}
optionally with deterministic samples on the projected $3\sigma$ ellipse to constrain the
footprint rather than the center.  \toshow{The tested nearest-ellipse hinge variant did not
reduce violations materially and is retired at that formulation; \cref{eq:soft-support} is a
design to freeze before use, not an evaluated method.  An independently measurable
surface-occupancy test for hull-interior floaters (depth reference, cross-view occupancy, or
geometry ground truth) remains an open obligation (O-12) --- ``no physical free-floaters''
must not be claimed from \cref{eq:containment} alone.}

\subsection{What the earlier RGB-backed schedules taught us}
\label{sec:carrier-history}

Two negative development results shaped this design and belong in the record.  (i) A short
fixed-budget carrier schedule (30/40/30/60 updates with clone waves, followed by a 40-update
recovery) \emph{preserved} its carriers ($94.6\%$ survival) yet finished $1.57$~dB below
immediate standard refinement; removing covariance repair \emph{improved} the endpoint by
$1.27$~dB even though the repair halved its own local residual ($-50.2\%$) --- local repair
objectives and downstream quality can dissociate when correspondences are
wrong.\evidence{ara/logic/claims.md C28}  (ii) The instrumented all-view maturation run
completed but violated its convergence-between-stages contract: three clone recoveries and
the higher-SH phase exhausted their 15k-update caps without reaching the plateau
criterion.\evidence{ara/logic/claims.md C29}  Both runs consumed RGB during optimization.
The compact-only sequence above replaces them: paired controls, predeclared budgets, no
topology growth, and a hard no-image boundary.  The dissociation lesson survives as a design
rule: \emph{every} local repair stage must justify itself by the matched no-repair ablation
on downstream risk (which the corrected repair now does, \cref{sec:cov-repair}), never by
its own residual.

\begin{figure}[t]
  \centering
  \figplaceholder{5.5cm}{%
    \textbf{Figure 9 --- carrier maturation stages (visual).}  One row per stage (Beam raw,
    after covariance repair, after phase 1, after prune, after phase 2): rendered front view
    + one side view, with carrier count and $\Jq$ annotated.  Purpose: make each accepted
    stage's contribution visible; the covariance-repair panel should visibly de-elongate
    beams.  Data: representative-root artifacts of the 2026-07-28 compact-only runs
    (\repopath{runs/}; re-render with the repository viewer).}
  \caption{\redcaption{Stage-by-stage carrier maturation renders.  To be produced from
    existing development artifacts.}}
  \label{fig:carrier-stages}
\end{figure}

\section{Experiments: Part II}
\label{sec:exp2}

\subsection{Design: Beam value cannot be inferred from a Beam-only result}
\label{sec:exp2-design}

The Beam hypothesis (\cref{claim:beam}) is tested only by a matched three-way initialization
comparison at exact count and compute, all downstream settings identical:
\begin{itemize}[leftmargin=1.6em,itemsep=1pt]
  \item \textbf{Beam Fusion + carriers} (treatment, \cref{sec:beam-fusion,sec:carrier});
  \item \textbf{balanced compact carve} (strong image-free structural control,
        \cref{sec:exp1-baselines}) --- same input information, no tomographic fusion;
  \item \textbf{calibration-bounded random} (lower bound) --- uniform in the camera-derived
        bounding volume.
\end{itemize}
All three feed the identical compact optimizer and evaluation samples.  Reported:
initialization quality, fixed intermediate checkpoints, risk-curve area (time-to-quality),
selected endpoint, paired seed wins, and both development scenes plus
\toshow{fresh confirmatory scenes}.  Beam construction time and peak memory are accounted
separately --- the initializer must pay for itself end-to-end, not hide its cost in
``preprocessing''.  A converged tie after an initialization win means Beam improves startup
but not final capacity; the text will say exactly that if it happens.

\begin{table}[t]
  \centering\small
  \begin{tabular}{@{}lccccc@{}}
    \toprule
    Initialization & init $\Jq$ & AUC($\Jq$) & final $\Jq$ & held-out PSNR & seed wins \\
    \midrule
    Beam + carriers          & \tbd & \tbd & \tbd & \tbd & \tbd \\
    Balanced compact carve   & \tbd & \tbd & \tbd & \tbd & \tbd \\
    Bounded random           & \tbd & \tbd & \tbd & \tbd & \tbd \\
    \bottomrule
  \end{tabular}
  \caption{\redcaption{Matched three-way initialization comparison (per scene; paired seeds;
    Beam build time and peak memory reported in a companion column).  Entirely pending ---
    this is the claim-deciding table of Part~II.}}
  \label{tab:beam-value}
\end{table}

\begin{figure}[t]
  \centering
  \figplaceholder{5.5cm}{%
    \textbf{Figure 10 --- time-to-quality curves.}  Validation $\Jq$ (y, log) versus
    optimizer updates (x) for the three initializations of \cref{tab:beam-value}; seeds as
    thin lines, medians bold; annotate initialization points and the fixed checkpoints.  A
    second panel with held-out PSNR versus wall-clock (including Beam construction time as
    an x-offset for the Beam arm) makes the ``pays for itself'' question visible.  Data: the
    matched comparison runs.}
  \caption{\redcaption{Risk trajectories for Beam versus carve versus random initialization.
    To be produced.}}
  \label{fig:time-to-quality}
\end{figure}

\subsection{Ablations}
\label{sec:exp2-ablations}

\begin{table}[t]
  \centering\footnotesize
  \setlength{\tabcolsep}{4pt}
  \begin{tabular}{@{}lp{7.1cm}l@{}}
    \toprule
    Ablation & Development verdict (single scene) & Confirmatory \\
    \midrule
    No corrected covariance repair
      & repair $-63.2\%$ $\Jq$, 3/3 wins \evidence{C30}
      & \tbd \\
    Legacy (one-sided, undilated) repair
      & $+51.6\%$ $\Jq$: harmful, retired \evidence{C30}
      & --- \\
    Legacy opacity / appearance repair
      & $+9.8\%$ / $+0.12\%$: retired \evidence{C30}
      & --- \\
    One compact phase vs.\ two
      & second phase $\to 0.717\times$ $\Jq$, 3/3 \evidence{C30}
      & \tbd \\
    Phase-1 means frozen
      & fails $2\%$ non-inferiority \evidence{C30}
      & \tbd \\
    Means-only optimization
      & $4.18\times$ $\Jq$: rejected \evidence{C30}
      & --- \\
    Without containment prune
      & prune costs $2.7\%$ $\Jq$, zero violations \evidence{C30, C31}
      & \tbd \\
    Prune before vs.\ after phase 2
      & before selected \evidence{C30}
      & \tbd \\
    Phase-2 means trainable vs.\ frozen
      & frozen passes ($1.015\times$ $\Jq$) \evidence{C30}
      & \tbd \\
    Clone (incl.\ density-preserving)
      & fails $5\%$ gate: no growth \evidence{C30}
      & --- \\
    SH0 vs.\ incremental SH3
      & $2.3$--$3.0\% < 5\%$ floor: SH0 \evidence{C30}
      & \tbd \\
    Soft support (\cref{eq:soft-support})
      & hinge variant immaterial, retired
      & \toshow{freeze then test} \\
    \bottomrule
  \end{tabular}
  \caption{Carrier-stage ablation grid.  Development verdicts are frozen paired-root results
    on one outcome-exposed scene (19 fitting views, compact validation; evidence keys refer
    to \repopath{ara/logic/claims.md}).  \toshow{The confirmatory column --- fresh scenes,
    held-out cameras, $\ge 3$ paired seeds, preregistered gates --- is entirely pending;
    rejected legacy stages re-enter only as negative controls.}}
  \label{tab:ablations}
\end{table}

The development column of \cref{tab:ablations} closes the \emph{implementation-policy}
question on one scene: which stages the carrier pipeline runs, in which order.  It does not
close the \emph{value} question (\cref{tab:beam-value}), which is a comparison against
non-Beam arms, nor any cross-scene generalization.

\subsection{Diagnostics to report}
\label{sec:exp2-diagnostics}

Per Beam run: contributor-count histogram, per-view uniqueness, gate survival rates,
covariance eigenvalue spectra, lineage hashes; per maturation: carrier survival, per-family
optimizer motion, covariance residual distribution (\cref{eq:cov-residual}), containment
violation counts before/after prune, and the fixed-topology / frozen-mean containment
recheck.  These are the quantities that make a false-association failure (the designed risk
of CI fusion, \cref{sec:ci-fusion}) visible instead of silent.

\begin{figure}[t]
  \centering
  \figplaceholder{5cm}{%
    \textbf{Figure 11 --- Beam diagnostics panel.}  (a) contributor-count histogram; (b)
    carrier survival across stages (bar per stage); (c) distribution of covariance residuals
    \cref{eq:cov-residual} before/after repair; (d) containment violations per view before
    prune.  Data: diagnostics dictionaries already emitted by the implementation on the
    development scene (\repopath{runs/}).}
  \caption{\redcaption{Beam Fusion and carrier diagnostics.  To be produced from existing
    run artifacts.}}
  \label{fig:beam-diagnostics}
\end{figure}

\subsection{Decision rule for Part II}
\label{sec:exp2-decision}

Preregistered: Beam is claimed as a contribution only if it beats the balanced carve control
on the frozen primary endpoint (held-out quality at matched compute and capacity) with the
agreed seed-win rule, replicated on fresh scenes.  \toshow{The exact endpoint, materiality
margin, and seed-win rule must be frozen before the confirmatory runs (obligation O-9).}
A tie or loss is reported in full as a bounded result --- the tomographic construction,
observability result, and lineage machinery remain contributions of record, but not
quality claims --- and, per \cref{sec:decision-rule}, has no bearing on Part~I.


\section{Limitations}
\label{sec:limitations}

\paragraph{Claim boundary.}
The VRAM claim covers post-fit reconstruction only.  Capture-time fitting costs real compute
(disclosed, \toshow{table pending, O-2}) and requires access to the full-resolution images
once.  Scenarios that need repeated access to original pixels (relighting, re-editing,
re-fitting at higher budgets) must retain the source data somewhere --- the claim is about
the reconstruction working set, not about deleting archives.

\paragraph{Teacher ceiling.}
Compact supervision can at best reproduce what the captures preserved.  Fit quality bounds
reconstruction quality; high-frequency texture beyond the component budget is unrecoverable
downstream.  The parity question of Part~I is therefore jointly a question about Stage-1
budgets, and the rate--distortion analysis (\cref{fig:rd}) is part of the claim's support,
not decoration.

\paragraph{View-dependent appearance.}
Current captures store one color model per 2D component fitted from one photograph; the
carrier path trains SH0 and found incremental SH3 immaterial \emph{under compact
supervision on one scene} (\cref{tab:ablations}).  Strongly view-dependent scenes
(specular, glass) may expose this as the binding constraint.
\toshow{Either demonstrate SH growth under compact supervision on a specular scene or scope
the claim to approximately-diffuse captures (O-11).}

\paragraph{Geometry guarantees.}
The containment invariant is a visual-hull center rule (\cref{sec:containment}): it cannot
detect hull-interior floaters and does not prove surface occupancy.  ``No free-floating
Gaussians'' in the physical sense requires an independent occupancy test
(\toshow{O-12}).

\paragraph{Gauge non-identifiability.}
Amplitude/color gauge freedom in the captures (\cref{sec:field-model}) means some 2D
quantities have no physical 3D reading; every stage that consumes raw amplitudes must be
gauge-audited.  Two former repair stages failed exactly this audit and are retired.

\paragraph{Association risk in Beam Fusion.}
CI fusion is conservative by construction, so Beam's principal failure mode is false
contributor association across views (repetitive texture, thin structures).  The gates of
\cref{eq:gates} are analytic but heuristic; the diagnostics of \cref{sec:exp2-diagnostics}
make failures visible, and the matched carve control quantifies what the tomographic step
adds.

\paragraph{Evidence maturity.}
All quantitative statements outside the mechanism results are single-scene, outcome-exposed
development evidence, and are labelled as such where they appear.  The confirmatory program
(\cref{sec:obligations}) --- fresh scenes, held-out cameras, paired seeds, preregistered
gates, controlled resource measurement, standard benchmarks, and independent audit --- is
the difference between this draft and a submittable paper.

\paragraph{Static scenes.}
The pipeline is static by design.  The carrier concept extends naturally toward dynamic
capture (per-frame 2D Gaussian video $\to$ space--time Beam Fusion $\to$ 4D carriers
$\to$ standard 4D Gaussian splatting); this paper intentionally establishes the static
foundation and defers temporal modeling.

\section{Conclusion}
\label{sec:conclusion}

We described a reconstruction pipeline whose supervision is a compact, analytic, frozen 2D
Gaussian field per view instead of the source photographs, and we split its evaluation into
two firewalled claims.  Part~I --- the systems claim --- argues that after a one-time
per-view fit, standard 3DGS optimization never needs the images again: supervision becomes a
stream of point queries against kilobyte-scale fields, and the dataset-scale RGB working set
disappears from the reconstruction process.  The mechanism is implemented end-to-end with a
verified point-rendering parity anchor, an unbiased sampled objective, and a hard structural
no-image boundary; \toshow{the claim itself awaits its controlled experiment --- held-out
parity plus measured peak-VRAM reduction under the preregistered protocol of
\cref{sec:vram-protocol}}.  Part~II --- Beam Fusion --- shows that the same captures support
a closed-form tomographic 3D initialization with a principled fusion rule, an observability
theorem fixing the minimum view count, complete contributor lineage, and a compact-only
carrier maturation whose every stage is justified by paired ablation; \toshow{whether the
initialization is \emph{worth having} against the matched no-Beam control is an open,
preregistered question, and a negative answer will be reported as such without touching
Part~I}.

The broader bet is that capture representations --- not model compression and not optimizer
machinery --- are the under-explored axis of scaling 3D reconstruction: once images are
replaced by queryable analytic fields, memory decouples from view count, supervision
decouples from resolution, and the same construction extends to time.  If the confirmatory
program supports Part~I, dome-scale and video-scale captures stop being a memory problem;
if it refutes it, the negative result will delimit exactly how much information per-view
2D Gaussian fits preserve --- either way, the decision rule was written down first.


\bibliographystyle{plainnat}
\bibliography{references}

\appendix
\section{Evidence Obligations (Tracking Table)}
\label{sec:obligations}

Every red item in the draft, in one place.  A row turns black only when its artifact exists,
passed independent audit, and has a claim-ledger row in \repopath{ara/logic/claims.md}
(repository evidence discipline).  This table is the draft's definition of done.

\begin{table}[h]
  \centering\small
  \begin{tabular}{@{}llp{6.6cm}l@{}}
    \toprule
    ID & Where & Obligation & Status \\
    \midrule
    O-1  & \cref{sec:protocol}
         & NVML sampler + fresh-process resource harness shared by compact and RGB arms
         & \toshow{open} \\
    O-2  & \cref{sec:claim-boundary}
         & Upstream Stage-1 cost disclosure table (fit time, VRAM, storage per view)
         & \toshow{open} \\
    O-3  & \cref{sec:decision-rule}
         & Preregistration: parity floor, VRAM materiality rule, multiplicity policy,
           frozen Arm-1 optimizer/topology/stopping rule
         & \toshow{open} \\
    O-4  & \cref{sec:exp-setup}
         & At least one outcome-unseen calibrated scene with untouched test cameras
         & \toshow{open} \\
    O-5  & \cref{sec:exp-setup}
         & Standard benchmarks (Mip-NeRF~360, Tanks and Temples) through Stage 1 + both arms
         & \toshow{open} \\
    O-6  & \cref{sec:exp1-results}
         & Main table + headline VRAM/views figure + qualitative parity grid
           (\cref{tab:main-results}, \cref{fig:vram-views}, \cref{fig:parity-grid})
         & \toshow{open} \\
    O-7  & \cref{sec:arms}
         & Real COLMAP sparse reconstructions for the Original-3DGS baselines
         & \toshow{open} \\
    O-8  & \cref{sec:exp2-design}
         & Matched Beam vs.\ carve vs.\ random comparison with separate Beam cost
           accounting (\cref{tab:beam-value}, \cref{fig:time-to-quality})
         & \toshow{open} \\
    O-9  & \cref{sec:exp2-decision}
         & Frozen Part-II endpoint, materiality margin, and seed-win rule
         & \toshow{open} \\
    O-10 & \cref{sec:exp2-ablations}
         & Confirmatory ablation column: fresh scenes, held-out cameras, $\ge 3$ paired
           seeds
         & \toshow{open} \\
    O-11 & \cref{sec:limitations}
         & View-dependent appearance: SH growth under compact supervision on a specular
           scene, or scope reduction
         & \toshow{open} \\
    O-12 & \cref{sec:containment}
         & Independent surface-occupancy / hull-interior floater test
         & \toshow{open} \\
    O-13 & \cref{sec:related}
         & Novelty ledger N1--N5: dated verdicts with citations; verify every bib entry
         & \toshow{open} \\
    O-14 & \cref{sec:point-rasterizer}
         & CUDA point-rasterizer parity (or explicit scope-out) + throughput numbers
         & \toshow{open} \\
    O-15 & \cref{sec:fitting}
         & Stage-1 statistics table + frozen fitter configuration; rate--distortion curve
           (\cref{fig:rd})
         & \toshow{open} \\
    \bottomrule
  \end{tabular}
  \caption{Open evidence obligations.  IDs are referenced from the running text.}
  \label{tab:obligations}
\end{table}

\section{Figure Production List}
\label{sec:figure-list}

Consolidated from the red figure slots (specifications in the respective captions):
\begin{enumerate}[leftmargin=1.8em,itemsep=1pt]
  \item \cref{fig:teaser} --- pipeline overview / claim boundary (drawn; thumbnails from
        existing captures);
  \item \cref{fig:capture} --- capture visualization: photo, field render, error, ellipses
        (from checked-in captures; producible now);
  \item \cref{fig:rd} --- capture rate--distortion vs.\ codecs (needs Stage-1 sweep);
  \item \cref{fig:sample-budget} --- quality vs.\ sample budget (needs frozen Arm-1
        schedule);
  \item \cref{fig:vram-views} --- \textbf{headline}: peak VRAM vs.\ view count (needs O-1,
        O-3, O-4);
  \item \cref{fig:parity-grid} --- qualitative held-out parity grid (needs O-6);
  \item \cref{fig:beam-geometry} --- Beam geometry schematic (drawable now, TikZ);
  \item \cref{fig:init-visual} --- initializer point-set comparison, front + side views
        (producible now from development artifacts);
  \item \cref{fig:carrier-stages} --- stage-by-stage carrier renders (producible now from
        2026-07-28 run artifacts);
  \item \cref{fig:time-to-quality} --- risk trajectories Beam/carve/random (needs O-8);
  \item \cref{fig:beam-diagnostics} --- Beam diagnostics panel (producible now from run
        diagnostics).
\end{enumerate}

\section{Artifact and Evidence Index (Draft-Only)}
\label{sec:artifact-index}

Black statements in this draft bind to the repository's audited evidence ledger.  Key
artifacts: \repopath{ara/logic/claims.md} (claim rows C15, C18--C20, C23, C24, C28--C31
cited in the text), \repopath{docs/ADR-002-carrier-refinement.md} (carrier policy and math
disposition), \repopath{docs/THREE_ARM_EXPERIMENT_PROGRAM.md} (arm design),
\repopath{docs/PAPER_PLAN_beam_fusion.md} (claim discipline and status),
\repopath{benchmarks/results/} (sealed results and independent audits).  This appendix is
removed from any submitted version.


\end{document}
